\begin{abstract}
How many points of monetary policy is a movement in financial conditions worth? I compare two answers. The first weights asset prices by how they respond to monetary policy news in announcement windows. The second weights them by their effects on GDP in the FRB/US model. Once the two are put on a common scale they agree on the combination of asset prices that follows a policy surprise, and the scale itself depends on how long the movements are assumed to last. For other combinations they differ. In point estimates the model gives the dollar more weight, relative to equities, than the event study does, but the size of the difference depends on how policy news is measured and the data do not establish it for each measure. The difference between the two readings is roughly 0.2 to 0.4 of the average market move and is about as large as the difference between two reasonable estimates of the market weights. In 1998--2001 the two readings of the same asset movements have opposite signs in the point estimates, both small; in 2022--24 both are about one point of tightening. Realized unemployment and aggregate GDP growth do not separate the readings, and the tests have little power. A reading of financial conditions in policy units is therefore identified up to a band, which the paper measures.
\end{abstract}

\noindent\textbf{JEL Classification:} E43, E44, E52, C32

\noindent\textbf{Keywords:} financial conditions index, monetary policy, FRB/US, event study, neutral interest rate

\pagestyle{plain}

\newpage
\pagenumbering{arabic}

% ============================================================

